% BHU PYQ -- Transcribed & Corrected
% Paper: POLMD-31 / POLMD31: Development of Indian Constitution
% Exam: Under Graduate Semester III Examination, 2025-26 (NEP)
% Category: Multidisciplinary Course (MD)
% Department: Political Science

\documentclass[11pt,a4paper]{article}
\usepackage[a4paper,top=2.5cm,bottom=2.5cm,left=2.8cm,right=2.8cm,headheight=14pt]{geometry}
\usepackage{amsmath,amssymb}
\usepackage[T1]{fontenc}
\usepackage[utf8]{inputenc}
\usepackage{lmodern,microtype,fancyhdr,enumitem,hyperref,lastpage,parskip}

\hypersetup{
    colorlinks=true,
    urlcolor=black,
    linkcolor=black,
    pdftitle={POLMD31: Development of Indian Constitution -- BHU Sem III 2025-26 (NEP MD)}
}

\pagestyle{fancy}
\fancyhf{}
\renewcommand{\headrulewidth}{0.4pt}
\renewcommand{\footrulewidth}{0.4pt}
\fancyhead[L]{\small   \textit{Banaras Hindu University}}
\fancyhead[R]{\small   \textit{POLMD-31: Development of Indian Constitution (MD)}}
\fancyfoot[C]{\small Page  \thepage\ of \pageref{LastPage}}


\newlist{parts}{enumerate}{2}
\setlist[parts,1]{label=(\alph*),leftmargin=*,itemsep=4pt,topsep=2pt}
\setlist[parts,2]{label=(\roman*),leftmargin=*,itemsep=2pt,topsep=2pt}

\newcommand{\pts}[1]{\hfill   \textit{[#1]}}


\begin{document}


\begin{center}
  {\large\bfseries BANARAS HINDU UNIVERSITY}\\[2pt]
  {\normalsize Under Graduate --- Semester III Examination, 2025-26 (NEP)}\\[6pt]
  \rule{0.8\linewidth}{0.5pt}\\[6pt]
  {\large\bfseries POLITICAL SCIENCE}\\[4pt]
  {\large\bfseries Paper Code: POLMD-31 : Development of Indian Constitution}\\[4pt]
  {\normalsize Time: 3 Hours\hfill Maximum Marks: 70}
\end{center}

\rule{\linewidth}{0.4pt}

\smallskip



\noindent   \textit{Instructions: Answer questions from Sections A, B, and C as per the instructions given in each section.}

\bigskip

\centerline{   \textbf{SECTION A}}
\smallskip



\noindent   \textit{Note: Answer each question within 50 words. Each question carries 2 marks.}

\bigskip




\noindent   \textbf{Question 1.} Short Concepts:\pts{$5  \times 2 = 10$}

\begin{parts}
  \item Pitt's India Act 1784\pts{2}
  \item Drafting Committee of the Indian Constitution\pts{2}
  \item Provincial Autonomy\pts{2}
  \item Constituent Assembly of India\pts{2}
  \item Interim Government of 1946\pts{2}
\end{parts}

\bigskip
\rule{\linewidth}{0.2pt}
\medskip

\centerline{   \textbf{SECTION B}}
\smallskip



\noindent   \textit{Note: Answer each question within 250 words. Each question carries 10 marks.}

\bigskip




\noindent   \textbf{Question 2.} Write the main provisions of the Regulating Act of 1773.\pts{10}

\centerline{   \textbf{OR}}




\noindent Write a detailed note on the features and impact of the Indian Councils Act of 1861.\pts{10}

\medskip\rule{\linewidth}{0.1pt}\medskip




\noindent   \textbf{Question 3.} Write a note on the key characteristics of the Indian Councils Act of 1909 (Morley-Minto Reforms).\pts{10}

\centerline{   \textbf{OR}}




\noindent Write a short note on the dyarchy system introduced in Indian provinces by the Government of India Act 1919.\pts{10}

\medskip\rule{\linewidth}{0.1pt}\medskip




\noindent   \textbf{Question 4.} Write a short note on the major developments in India's Independence Movement during 1947.\pts{10}

\centerline{   \textbf{OR}}




\noindent Write a short note on the main provisions and significance of the Indian Independence Act 1947.\pts{10}

\bigskip
\rule{\linewidth}{0.2pt}
\medskip

\centerline{   \textbf{SECTION C}}
\smallskip



\noindent   \textit{Note: Answer all long essay questions. Each question carries 15 marks.}

\bigskip




\noindent   \textbf{Question 5.} Examine the historical evolution of the Indian Constitution by tracing the key British Acts passed before independence.\pts{15}

\centerline{   \textbf{OR}}




\noindent Write a comprehensive essay on the drafting process of the Indian Constitution, from the formation of the Constituent Assembly to its final adoption.\pts{15}

\medskip\rule{\linewidth}{0.1pt}\medskip




\noindent   \textbf{Question 6.} Write an essay on how the Government of India Act 1935 increased Indian participation in governance and established provincial autonomy.\pts{15}

\centerline{   \textbf{OR}}




\noindent Explain why the Government of India Act 1935 is considered the most critical constitutional milestone prior to Indian independence.\pts{15}

\vfill
\rule{\linewidth}{0.4pt}

\begin{center}\small
\href{https://bhu-exam.vercel.app/}{Click here to visit our website}\\[4pt]
\href{https://me-aryan.vercel.app/}{Click here to visit our contributor}
\end{center}

\end{document}
